\chapter{Calcium Pyrophosphate Deposition Disease}
\index{Calcium Pyrophosphate Deposition Disease}
\label{sec:Calcium Pyrophosphate Deposition Disease}

\section{Overview}
Calcium pyrophosphate deposition disease (CPPD, formerly pseudogout) is a crystal-induced arthropathy resulting from the deposition of calcium pyrophosphate dihydrate (CPP) crystals in articular cartilage (chondrocalcinosis) and synovium. The condition occurs predominantly in elderly individuals ($>60$ years) and is linked to hyperparathyroidism, hemochromatosis, hypomagnesemia, and hypothyroidism.
\section{Causes}
\section{Risk Factors}

\begin{itemize}[noitemsep]
  \item Genetic predisposition and family history
  \item Advancing age
  \item Lifestyle factors (diet, exercise, smoking, alcohol)
  \item Environmental exposures
  \item Underlying medical conditions
  \item Medication use
\end{itemize}
\section{Signs and Symptoms}

\subsection{Common symptoms}
\begin{itemize}[noitemsep]
  \item You may experience acute monoarthritis or oligoarthritis resembling gout, most commonly affecting the knee, wrist, or symphysis pubis, with swelling, warmth, severe pain, and redness of the skin. 
  \end{itemize}

\subsection{Emergency warning signs}
\begin{itemize}[noitemsep]
  \item Severe or worsening symptoms of calcium pyrophosphate deposition disease require immediate medical evaluation
\end{itemize}
\section{Progression and Stages}
The condition may progress through different stages of severity over time. Early stages often involve mild or subtle symptoms that may be overlooked. Moderate stages involve more noticeable symptoms affecting daily function. Advanced stages may involve significant impairment and complications.
\section{Complications}
\begin{itemize}[noitemsep]
  \item Disease progression leading to worsening symptoms
  \item Secondary infections or injuries
  \item Side effects of treatment
  \item Reduced quality of life
  \item Emotional and psychological impact
\end{itemize}
\section{Diagnosis}
\begin{itemize}[noitemsep]
  \item Your doctor confirms the diagnosis with synovial fluid removal by suction showing weakly positively birefringent, rhomboid-shaped CPP crystals under polarized light microscopy
  \item Radiographs demonstrate chondrocalcinosis (linear calcium buildup of articular cartilage). Treatment for acute attacks includes intra-articular corticosteroid injection, oral NSAIDs, or colchicine.
  \item Comprehensive medical history and physical examination
  \item Laboratory tests to assess organ function
  \item Imaging studies to visualize affected structures
  \item Specialized testing depending on the affected system
  \item Consultation with relevant specialists
\end{itemize}
\section{Prevention}
\begin{itemize}[noitemsep]
  \item Maintain a healthy lifestyle with balanced diet and exercise
  \item Avoid known risk factors
  \item Attend regular medical check-ups for early detection
  \item Follow recommended screening guidelines
\end{itemize}
\section{Treatment Options}
Medications to manage symptoms and address underlying mechanisms Lifestyle modifications and self-care measures Physical therapy and rehabilitation Surgical intervention when indicated Regular monitoring and follow-up care Patient education and support services

\begin{itemize}[noitemsep]
  \item Medications to manage symptoms and address underlying mechanisms
  \item Lifestyle modifications and self-care measures
  \item Physical therapy and rehabilitation
  \item Surgical intervention when indicated
  \item Regular monitoring and follow-up care
\end{itemize}
\section{Living with It}
\begin{itemize}[noitemsep]
  \item Follow your treatment plan as prescribed
  \item Maintain regular follow-up with your healthcare provider
  \item Make necessary lifestyle adjustments
  \item Seek support from family, friends, or support groups
  \item Stay informed about your condition
\end{itemize}
\section{Outlook}
The outlook is good for acute control. The outlook varies depending on the specific condition, severity at diagnosis, and response to treatment. Early intervention generally leads to better outcomes. Many conditions can be effectively managed with appropriate treatment. Regular follow-up care is important for monitoring and treatment adjustment.
